%------------------------------------------------------------------------------%
\documentclass[a4paper,12pt,fleqn]{article}

\usepackage{gaal}

%\ShowAnswers

%------------------------------------------------------------------------------%
\begin{document}

\makeHeader{GAAL}{Prova 1 -- Turma 7}{26/08/2026}

% 01
%------------------------------------------------------------------------------%
\exercise{25\%}
Dados
\[
  A =
  \begin{bmatrix}
    1 & 2 & -1 \\
    0 & 3 & 2  \\
    1 & 1 & 3
  \end{bmatrix}
\qquad
  \mathbf{x} =
  \begin{bmatrix}
    x \\ y \\ z
  \end{bmatrix}
\]
calcule \(\mathbf{x}^t A \mathbf{x}\)

\begin{answer}
  \begin{align*}
    \mathbf{x}^t A
    & =
    \begin{bmatrix}
      x \\ y \\ z
    \end{bmatrix}^t
    \begin{bmatrix}
      1 & 2 & -1 \\
      0 & 3 & 2  \\
      1 & 1 & 3
    \end{bmatrix}
    \\[2mm]
    & =
    \begin{bmatrix}
      x & y & z
    \end{bmatrix}
    \begin{bmatrix}
      1 & 2 & -1 \\
      0 & 3 & 2  \\
      1 & 1 & 3
    \end{bmatrix}
    \\[2mm]
    & =
    \begin{bmatrix}
      x \cdot 1 + y \cdot 0 + z \cdot 1 &
      x \cdot 2 + y \cdot 3 + z \cdot 2 &
      x \cdot 1 + y \cdot 1 + z \cdot 3
    \end{bmatrix}
    \\[2mm]
    & =
    \begin{bmatrix}
      x + z &
      2x + 3y + 2z &
      x + y + 3z
    \end{bmatrix}
  \end{align*}
  \begin{align*}
    \mathbf{x}^t A \mathbf{x}
    & =
    \begin{bmatrix}
      x + z &
      2x + 3y + 2z &
      x + y + 3z
    \end{bmatrix}
    \begin{bmatrix}
      x \\ y \\ z
    \end{bmatrix}
    \\[2mm]
    & =
    \begin{bmatrix}
      (x + z)x +
      (2x + 3y + 2z)y +
      {x + y + 3z}z
    \end{bmatrix}
    \\[2mm]
    & =
      x^2 + xz +
      2xy + 3y^2 + 2yz +
      xz + yz + 3z^2
    \\[2mm]
    & =
      x^2 + 3y^2 + 3z^2 + 2xy + 2xz + + 3yz
  \end{align*}
\end{answer}

% 02
%------------------------------------------------------------------------------%
\exercise{25\%}
Resolva o sistema linear por Gauss--Jordan
\[
  \systeme{
     x + 2y -  z = 0,
    2x + 4y - 2z = 0,
     x +  y +  z = 0
  }
\]

\begin{answer}
Vamos escalonar a matriz do sistema
\[
  \begin{bmatrix}
    1 & 2 & -1 \\
    2 & 4 & -2 \\
    1 & 1 &  1
  \end{bmatrix}
\]
Eliminando os elementos abaixo do primeiro pivô
\begin{gather*}
  L_2 \leftarrow L_2 - 2 L_1 \\
  L_3 \leftarrow L_3 - L_1
\end{gather*}
\[
\begin{bmatrix}
  1 & 2  & -1 \\
  0 & 0  & 0  \\
  0 & -1 & 2
\end{bmatrix}
\]
Trocando $L_2$ e $L_3$
\[
\begin{bmatrix}
  1 & 2  & -1 \\
  0 & -1 & 2  \\
  0 & 0  & 0
\end{bmatrix}
\]
Igualando o segundo pivô a 1
\[
  L_2 \leftarrow - L_2
\]
\[
\begin{bmatrix}
  1 & 2 & -1 \\
  0 & 1 & -2 \\
  0 & 0 & 0
\end{bmatrix}
\]
Eliminando os elementos acima do segundo pivô
\[
  L_1 \leftarrow L_1 - 2 L_2
\]
\[
\begin{bmatrix}
  1 & 0 &  3 \\
  0 & 1 & -2 \\
  0 & 0 &  0
\end{bmatrix}
\]
A variável $z$ é livre.

Fazemos $z=t$
\[
  x= -3t
  \qquad
  y = 2t
\]
A solução é
\[
  \mathbf{x}
  =
  t\begin{bmatrix}
    -3 \\
    2  \\
    1
  \end{bmatrix}
  \qquad t \in \R
\]
\end{answer}

% 03
%------------------------------------------------------------------------------%
\exercise{25\%}
Utilize o método de Gauss--Jordan para inverter a matriz
\[
  A =
  \begin{bmatrix}
    1 & 1 & 0 \\
    0 & 1 & 1 \\
    1 & 0 & 1
  \end{bmatrix}
\]

\begin{answer}
Matriz aumentada
\[
  \begin{amatrix}[3]{3}
    1 & 1 & 0 & 1 & 0 & 0 \\
    0 & 1 & 1 & 0 & 1 & 0 \\
    1 & 0 & 1 & 0 & 0 & 1
  \end{amatrix}
\]
Aplicamos
\[
  L_3 \leftarrow L_3 - L_1
\]
\[
  \begin{amatrix}[3]{3}
    1 & 1  & 0 & 1  & 0 & 0 \\
    0 & 1  & 1 & 0  & 1 & 0 \\
    0 & -1 & 1 & -1 & 0 & 1
  \end{amatrix}
\]
Agora
\[
 L_3 \leftarrow L_3 + L_2
\]
Obtemos
\[
  \begin{amatrix}[3]{3}
    1 & 1 & 0 & 1  & 0 & 0 \\
    0 & 1 & 1 & 0  & 1 & 0 \\
    0 & 0 & 2 & -1 & 1 & 1
  \end{amatrix}
\]
Dividimos a terceira linha por $2$
\[
  L_3 \leftarrow \sfrac{1}{2} L_3
\]
Logo
\[
  \begin{amatrix}[3]{3}
    1 & 1 & 0 & 1            & 0           & 0           \\
    0 & 1 & 1 & 0            & 1           & 0           \\
    0 & 0 & 1 & -\sfrac{1}{2} & \sfrac{1}{2} & \sfrac{1}{2}
  \end{amatrix}
\]
Eliminamos o elemento acima do terceiro pivô
\[
  L_2 \leftarrow L_2 - L_3
\]
Então
\[
  \begin{amatrix}[3]{3}
    1 & 1 & 0 & 1            & 0           & 0            \\
    0 & 1 & 0 & \sfrac{1}{2}  & \sfrac{1}{2} & -\sfrac{1}{2} \\
    0 & 0 & 1 & -\sfrac{1}{2} & \sfrac{1}{2} & \sfrac{1}{2}
  \end{amatrix}
\]
Finalmente
\[
  L_1 \leftarrow L_1 - L_2
\]
Assim
\[
  \begin{amatrix}[3]{3}
    1 & 0 & 0 & \sfrac{1}{2}  & -\sfrac{1}{2} & \sfrac{1}{2}  \\
    0 & 1 & 0 & \sfrac{1}{2}  & \sfrac{1}{2}  & -\sfrac{1}{2} \\
    0 & 0 & 1 & -\sfrac{1}{2} & \sfrac{1}{2}  & \sfrac{1}{2}
  \end{amatrix}
\]
Portanto
\[
  A^{-1} =
  \begin{bmatrix}
    \sfrac{1}{2}  & -\sfrac{1}{2} & \sfrac{1}{2}  \\
    \sfrac{1}{2}  & \sfrac{1}{2}  & -\sfrac{1}{2} \\
    -\sfrac{1}{2} & \sfrac{1}{2}  & \sfrac{1}{2}
  \end{bmatrix}
\]
\end{answer}

% 04
%------------------------------------------------------------------------------%
\exercise{25\%}
Utilize a inversa para resolver o sistema linear
\[
  \systeme{
  x + y = 7,
  y + z = 4,
  x + z = 5
  }
\]
\hint{Veja o exercício anterior}

\begin{answer}
A matriz do sistema é
\[
  A =
  \begin{bmatrix}
    1 & 1 & 0 \\
    0 & 1 & 1 \\
    1 & 0 & 1
  \end{bmatrix}
\]
No exercício anterior calculamos
\[
  A^{-1} =
  \begin{bmatrix}
    \sfrac{1}{2}  & -\sfrac{1}{2} & \sfrac{1}{2}  \\
    \sfrac{1}{2}  & \sfrac{1}{2}  & -\sfrac{1}{2} \\
    -\sfrac{1}{2} & \sfrac{1}{2}  & \sfrac{1}{2}
  \end{bmatrix}
\]
A solução do sistema é
\begin{align*}
  \mathbf{x}
  & = A^{-1} \mathbf{b}
  \\[2mm]
  & =
    \begin{bmatrix}
      \sfrac{1}{2}  & -\sfrac{1}{2} & \sfrac{1}{2}  \\
      \sfrac{1}{2}  & \sfrac{1}{2}  & -\sfrac{1}{2} \\
      -\sfrac{1}{2} & \sfrac{1}{2}  & \sfrac{1}{2}
    \end{bmatrix}
    \begin{bmatrix}
      7 \\
      4 \\
      5 \\
    \end{bmatrix}
  \\[2mm]
  & =
    \begin{bmatrix}
       \dfrac{1}{2}\, 7 - \dfrac{1}{2}\, 4 + \dfrac{1}{2}\, 5 \\[3mm]
       \dfrac{1}{2}\, 7 + \dfrac{1}{2}\, 4 - \dfrac{1}{2}\, 5 \\[3mm]
      -\dfrac{1}{2}\, 7 + \dfrac{1}{2}\, 4 + \dfrac{1}{2}\, 5
    \end{bmatrix}
  \\[2mm]
  & =
    \begin{bmatrix}
      \dfrac{ 7 - 4 + 5}{2} \\[3mm]
      \dfrac{ 7 + 4 - 5}{2} \\[3mm]
      \dfrac{-7 + 4 + 5}{2}
    \end{bmatrix}
  \\[2mm]
  & =
    \begin{bmatrix}
      \dfrac{8}{2} \\[3mm]
      \dfrac{6}{2} \\[3mm]
      \dfrac{2}{2}
    \end{bmatrix}
  \\[2mm]
  & =
    \begin{bmatrix}
      4 \\[3mm]
      3 \\[3mm]
      2
    \end{bmatrix}
\end{align*}
\end{answer}

%------------------------------------------------------------------------------%
\end{document}
%------------------------------------------------------------------------------%
